\subsection{Mobile Telegraphs}

\subsubsection{Описание решения}

Задача легко решается алгоритмом Дейкстры, если преобразовать входные данные в ориентированный взвешенный граф \( G=(V,E) \), где \( v\in V \) --- бойцы на поле с телеграфами, а \( (u,v)\in E \) --- связь между телеграфами \( u \) и \( v \). 

По условию задачи, две вершины в графе связаны, если их строковые номера различаются:
\begin{itemize}
  \item одной цифрой;
  \item двумя цифрами, которые можно поменять местами.
\end{itemize}

При этом вес ребра зависит от длины общего префикса и определяется массивом из входных данных.

Восстановление сертификата пути происходит за линейное время путём перехода по указателю на предыдущую вершину, которая участвовала в релаксации пути.

\subsubsection{Асимптотическая сложность}

\textit{Временная сложность решения} --- \( \mathcal{O}(N\cdot\log N+M\cdot\log N) \). Граф инициализируется с применением оптимизации на базе знания структуры входных строк: если проверять все пары строк и проверять условие, то решение быстро упрётся во временные рамки. Поэтому сортируем входной массив в лексикографическом порядке и ищем похожие строки при помощи бинарного поиска.

Также используется алгоритм Дейкстры с применением очереди с приоритетами, который проходится по всем достижимым рёбрам графа из стартовой вершины и определяет минимальных соседей за логарифмическое время.

\textit{Пространственная сложность решения} --- линейная \( \mathcal{O}(N) \). Все вершины материализованы в графе, а число рёбер для каждой вершины не более чем константно из-за специфики построения графа по строкам.

\subsubsection{Листинг кода}

\begin{lstlisting}
#include <algorithm>
#include <cstddef>
#include <iostream>
#include <limits>
#include <queue>
#include <stdexcept>
#include <vector>

namespace {
const std::size_t STRING_SIZE = 10;
const std::size_t INF = std::numeric_limits<std::size_t>::max();
}  // namespace

struct AdjacentMeta {
  bool are_adjacent = false;
  std::size_t common_prefix = 0;
};

struct Edge {
  std::size_t to_idx;
  std::size_t weight;
};

struct Vertex {
  std::size_t id;
  Vertex* prev = nullptr;
  std::vector<Edge> neighbours;
};

struct DijkstraResult {
  std::size_t dist;
  Vertex* end;
};

struct SearchNode {
  std::string str;
  std::size_t id;
  bool operator<(const SearchNode& other) const {
    return str < other.str;
  }
};

class Graph {
public:
  Graph(const std::vector<std::string>& strings, const std::vector<std::size_t>& costs)
      : vertices_(std::vector<Vertex>(strings.size() + 1)) {
    for (std::size_t i = 0; i < strings.size(); ++i) {
      vertices_[i].id = i;
    }

    std::vector<SearchNode> search_array(strings.size());
    for (std::size_t i = 0; i < strings.size(); ++i) {
      search_array[i] = {strings[i], i};
    }
    std::sort(search_array.begin(), search_array.end());

    auto find_id = [search_array](const std::string& target_str) -> std::size_t {
      SearchNode target{target_str, 0};
      auto it = std::lower_bound(search_array.begin(), search_array.end(), target);
      if (it != search_array.end() && it->str == target_str) {
        return it->id;
      }
      return INF;
    };

    for (std::size_t u = 0; u < strings.size(); ++u) {
      std::string s = strings[u];
      for (std::size_t i = 0; i < STRING_SIZE; ++i) {
        const char original_char = s[i];
        for (char d = '0'; d <= '9'; ++d) {
          if (d == original_char) {
            continue;
          }
          s[i] = d;
          std::size_t v = find_id(s);
          if (v != INF) {
            vertices_[u].neighbours.push_back({v, costs[i]});
          }
        }
        s[i] = original_char;
      }

      for (std::size_t i = 0; i < STRING_SIZE; ++i) {
        for (std::size_t j = i + 1; j < STRING_SIZE; ++j) {
          if (s[i] == s[j]) {
            continue;
          }

          std::swap(s[i], s[j]);
          std::size_t v = find_id(s);
          if (v != INF) {
            vertices_[u].neighbours.push_back({v, costs[i]});
          }
          std::swap(s[i], s[j]);
        }
      }
    }
  }

  static const AdjacentMeta AreAdjacent(std::string_view left, std::string_view right) {
    if (left.size() != right.size()) {
      throw std::invalid_argument("strings should have equal length");
    }
    std::size_t left_idx = 0;
    AdjacentMeta meta;
    while (left_idx < STRING_SIZE && left[left_idx] == right[left_idx]) {
      left_idx++;
      meta.common_prefix++;
    }
    std::size_t right_idx = left_idx + 1;
    while (right_idx < STRING_SIZE && left[right_idx] == right[right_idx]) {
      right_idx++;
    }
    std::size_t third_idx = right_idx + 1;
    while (third_idx < STRING_SIZE && left[third_idx] == right[third_idx]) {
      third_idx++;
    }
    meta.are_adjacent = (left_idx == STRING_SIZE) || (right_idx == STRING_SIZE) ||
                        (third_idx == STRING_SIZE && left[left_idx] == right[right_idx] &&
                         left[right_idx] == right[left_idx]);
    return meta;
  }

  DijkstraResult Dijkstra(std::size_t start_idx, std::size_t end_idx) {
    std::vector<std::size_t> dist(vertices_.size(), INF);
    std::priority_queue<
        std::pair<std::size_t, std::size_t>,
        std::vector<std::pair<std::size_t, std::size_t>>,
        std::greater<std::pair<std::size_t, std::size_t>>>
        pq;
    dist[start_idx] = 0;
    pq.emplace(0, start_idx);
    while (!pq.empty()) {
      auto [current_dist, u_idx] = pq.top();
      pq.pop();
      if (u_idx == end_idx) {
        return {dist[end_idx], &vertices_[end_idx]};
      }
      if (current_dist > dist[u_idx]) {
        continue;
      }
      for (const auto& edge : vertices_[u_idx].neighbours) {
        const std::size_t v_idx = edge.to_idx;
        if (dist[u_idx] + edge.weight < dist[v_idx]) {
          dist[v_idx] = dist[u_idx] + edge.weight;
          vertices_[v_idx].prev = &vertices_[u_idx];
          pq.emplace(dist[v_idx], v_idx);
        }
      }
    }
    return {dist[end_idx], &vertices_[end_idx]};
  }

private:
  std::vector<Vertex> vertices_;
};

int main() {
  std::size_t n_count = 0;
  std::cin >> n_count;
  std::vector<std::size_t> costs(STRING_SIZE);
  for (std::size_t i = 0; i < STRING_SIZE; i++) {
    std::cin >> costs[i];
  }
  std::vector<std::string> strings(n_count);
  for (std::size_t i = 0; i < n_count; i++) {
    std::cin >> strings[i];
  }
  Graph graph(strings, costs);
  const DijkstraResult result = graph.Dijkstra(0, n_count - 1);
  if (result.dist == INF) {
    std::cout << -1;
    return 0;
  }
  std::vector<std::size_t> path;
  const Vertex* curr = result.end;
  std::size_t count = 0;
  while (curr != nullptr) {
    path.push_back(curr->id);
    curr = curr->prev;
    count++;
  }
  std::reverse(path.begin(), path.end());
  std::cout << result.dist << '\n';
  std::cout << count << '\n';
  for (const std::size_t vertex_id : path) {
    std::cout << vertex_id + 1 << " ";
  }
  return 0;
}
\end{lstlisting}